% Ready-to-input fragment; no canonical files are modified by this artifact.
% Uses the existing q444:potential and q444:moving-family data and standard
% theorem/proposition/lemma environments. All new calculations are proved here.
\subsection{The full homogeneous gauge equations and the literal periodic orbit}

Retain signature $(+---)$, $\mathcal A_\mu\in\mathfrak{su}(3)$,
$\Phi=\Phi^*$, $\operatorname{tr}\Phi=0$, and $b\in\mathbb C^3$.
Write $g^2>0$ for the gauge coupling in this subsection and keep
\[
 s=\operatorname{tr}\Phi^2+2b^*b,\qquad
 a=m^2+\Gamma s,\qquad
 V=\frac{m^2}{2}s+\frac{\Gamma}{4}s^2-\kappa(\det\Phi-b^*\Phi b).
\]
Thus $m^2\geq0$, $\Gamma>0$, and $\kappa\in\mathbb R$ are unchanged.
For a matrix $Z$ write
$Z_{\mathrm{tf}}=Z-\operatorname{tr}(Z)I_3/3$. The full density is
\begin{equation}\label{q444:full-gauge-density}
 \mathcal L=\frac1{2g^2}\operatorname{tr}
 (\mathcal F_{\mu\nu}\mathcal F^{\mu\nu})
 +\frac12\operatorname{tr}(D_\mu\Phi D^\mu\Phi)
 +(D_\mu b)^*D^\mu b-V,
\end{equation}
where $\mathcal F_{\mu\nu}=\partial_\mu\mathcal A_\nu-
\partial_\nu\mathcal A_\mu+[\mathcal A_\mu,\mathcal A_\nu]$,
$D_\mu\Phi=\partial_\mu\Phi+[\mathcal A_\mu,\Phi]$, and
$D_\mu b=\partial_\mu b+\mathcal A_\mu b$.
We consider spatially homogeneous fields on a finite-volume flat
$\mathbb T^3$, or locally on $\mathbb R^3$. A nonzero constant energy
density is not asserted to have finite total energy on $\mathbb R^3$.

\begin{proposition}[Every homogeneous field equation]\label{q444:gauge-eom}
Set
\[
 E_i=\dot{\mathcal A}_i+[\mathcal A_0,\mathcal A_i],\qquad
 F_{ij}=[\mathcal A_i,\mathcal A_j],\qquad
 X=\dot\Phi+[\mathcal A_0,\Phi],\qquad v=\dot b+\mathcal A_0b,
\]
and define the lower-index matter currents
\[
 J_0=[\Phi,X]+bv^*-vb^*,\qquad
 J_i=[\Phi,[\mathcal A_i,\Phi]]
                -bb^*\mathcal A_i-\mathcal A_i bb^*.
\]
The complete homogeneous Euler--Lagrange system is
\begin{align}
 D_0X-\sum_i[\mathcal A_i,[\mathcal A_i,\Phi]]
       +a\Phi-\kappa(\operatorname{adj}\Phi-bb^*)_{\mathrm{tf}}&=0,
       \label{q444:gauge-phi}\\
 D_0v-\sum_i\mathcal A_i^2b+(aI_3+\kappa\Phi)b&=0,
       \label{q444:gauge-b}\\
 \frac2{g^2}\sum_i[\mathcal A_i,E_i]-(J_0)_{\mathrm{tf}}&=0,
       \label{q444:full-gauss}\\
 D_0E_i-\sum_j[\mathcal A_j,[\mathcal A_j,\mathcal A_i]]
       -\frac{g^2}{2}(J_i)_{\mathrm{tf}}&=0
       \quad(i=1,2,3). \label{q444:gauge-spatial}
\end{align}
Here $D_0X=\dot X+[\mathcal A_0,X]$,
$D_0v=\dot v+\mathcal A_0v$, and similarly for $E_i$.
\end{proposition}
\begin{proof}
First retain the homogeneous density without deleting $\mathcal A_0$:
\[
 \mathcal L_{\mathrm{hom}}
 =-\frac1{g^2}\sum_i\operatorname{tr}E_i^2
   +\frac1{g^2}\sum_{i<j}\operatorname{tr}F_{ij}^2
   +\frac12\operatorname{tr}X^2+v^*v
   -\frac12\sum_i\operatorname{tr}[\mathcal A_i,\Phi]^2
   -\sum_i|\mathcal A_i b|^2-V.
\]
Both signs in the gauge terms follow by raising indices in $(+---)$
and counting the two orders of each antisymmetric pair.
For trace-free Hermitian $\delta\Phi$ and arbitrary $\delta b$,
\[
 \delta V=
 \operatorname{tr}\!\left[
 \bigl(a\Phi-\kappa(\operatorname{adj}\Phi-bb^*)_{\mathrm{tf}}\bigr)
 \delta\Phi\right]
 +2\operatorname{Re}\bigl(\delta b^*(aI_3+\kappa\Phi)b\bigr).
\]
Covariant integration by parts therefore gives
$D_\mu D^\mu\Phi+a\Phi-
\kappa(\operatorname{adj}\Phi-bb^*)_{\mathrm{tf}}=0$ and
$D_\mu D^\mu b+(aI_3+\kappa\Phi)b=0$. Homogeneity gives the first
two displayed equations, including the minus signs in their spatial sums.

For the gauge variation, antisymmetry gives
\[
 \delta\mathcal L_{\mathrm{gauge}}
 =\frac2{g^2}\operatorname{tr}
       ((D_\mu\delta\mathcal A_\nu)\mathcal F^{\mu\nu}).
\]
After integration by parts its coefficient is
$-2D_\mu\mathcal F^{\mu\nu}/g^2$.
The scalar coefficient of $\delta\mathcal A_\nu$ is
\[
 J^\nu=[\Phi,D^\nu\Phi]
           +b(D^\nu b)^*-(D^\nu b)b^*.
\]
Indeed cyclicity gives
$\operatorname{tr}(D^\nu\Phi[\delta\mathcal A_\nu,\Phi])
=\operatorname{tr}([\Phi,D^\nu\Phi]\delta\mathcal A_\nu)$,
and anti-Hermiticity gives
$2\operatorname{Re}((D^\nu b)^*\delta\mathcal A_\nu b)
=\operatorname{tr}((b(D^\nu b)^*-(D^\nu b)b^*)\delta\mathcal A_\nu)$.
Variation over $\mathfrak{su}(3)$ consequently yields
\[
 D_\mu\mathcal F^{\mu\nu}=\frac{g^2}{2}(J^\nu)_{\mathrm{tf}}.
\]
For $\nu=0$, $\mathcal F^{i0}=E_i$ and $J^0=J_0$.
For $\nu=i$, $\mathcal F^{0i}=-E_i$,
$\mathcal F^{ji}=[\mathcal A_j,\mathcal A_i]$, and $J^i=-J_i$.
These substitutions give exactly \eqref{q444:full-gauss} and
\eqref{q444:gauge-spatial}. In particular the $\mathcal A_0$ equation
survives after choosing temporal gauge. This is the same variational
principle emphasized by D.~Tong, \emph{Lectures on Gauge Theory},
Section~2.2.1, equation~(2.27); the matter currents, signs, and coefficient
$g^2/2$ above have been derived in the present conventions.
\end{proof}

\begin{theorem}[No temporal-gauge lift retaining the literal orbit]
\label{q444:temporal-no-lift}
Keep exactly the moving-kernel scalar functions in
Theorem~\ref{q444:moving-family}, with $B,d>0$ and $x\neq0$.
There are no homogeneous spatial fields $\mathcal A_i(t)$ for which
these same scalar functions, together with $\mathcal A_0=0$, solve
the full system \eqref{q444:gauge-phi}--\eqref{q444:gauge-spatial}.
\end{theorem}
\begin{proof}
The scalar functions already solve the unrestricted scalar equations.
Subtracting them from \eqref{q444:gauge-phi}--\eqref{q444:gauge-b},
with $\mathcal A_0=0$, gives
\[
 \sum_i[\mathcal A_i,[\mathcal A_i,\Phi]]=0,\qquad
 \sum_i\mathcal A_i^2b=0.
\]
Taking the trace pairing with $\Phi$ and the Hermitian pairing with $b$
gives respectively
\[
 -\sum_i\operatorname{tr}[\mathcal A_i,\Phi]^2=0,\qquad
 -\sum_i|\mathcal A_i b|^2=0.
\]
Each $[\mathcal A_i,\Phi]$ is Hermitian, so both are sums of
nonnegative squared norms with the displayed common minus sign. Thus
$[\mathcal A_i,\Phi]=0$ and $\mathcal A_i b=0$ separately for every $i$.

To classify this stabilizer use the retained orthogonal matrix
$K=(v/\sqrt2,w_0/\sqrt6,u/\sqrt3)$ from the moving-family proof, and
put $r^2=x^2+d^2>0$. At time zero the data are
\[
 H=K^*\Phi_0K=\begin{pmatrix}0&0&0\\0&-x&d\\0&d&x\end{pmatrix},
 \qquad K^*b_0=Be_1,\qquad
 A=K^*\Omega K=
 \begin{pmatrix}0&-\sqrt a&0\\\sqrt a&0&0\\0&0&0\end{pmatrix}.
\]
An anti-Hermitian traceless matrix killing $Be_1$ has zero first row
and column and a lower $\mathfrak{su}(2)$ block. Commutation with
$\left(\begin{smallmatrix}-x&d\\d&x\end{smallmatrix}\right)$, whose
distinct eigenvalues are $\pm r$, forces that block to be its real
multiple times $i$. Therefore the stabilizer Lie algebra is
$\mathbb R\,i\Phi(t)$. The full stabilizer group is the circle
\[
 \left\{P_b+\cos\theta(I_3-P_b)
          +i\frac{\sin\theta}{r}\Phi(t):\theta\in\mathbb R\right\},
 \qquad P_b=\frac{b(t)b(t)^*}{B^2}.
\]
Indeed it fixes $b$, has eigenvalues $1,e^{i\theta},e^{-i\theta}$,
and exhausts the unitary commutant fixing the zero eigenspace pointwise.
Consequently
\[
 \mathcal A_i=i c_i(t)\Phi(t),\qquad c_i(t)\in\mathbb R.
\]
The original orbit has
\[
 [\Phi,\dot\Phi]=-r^2\Omega,\qquad
 b\dot b^*-\dot b b^*=-B^2\Omega,\qquad
 J_0=-(r^2+B^2)\Omega.
\]
For verification, $H^2=r^2\operatorname{diag}(0,1,1)$,
$HAH=0$, and hence $[H,[A,H]]=-r^2A$; the vector term
is $-B^2A$. Conjugation back by $O(t)K$, with the rotation commuting
with $\Omega$, gives the stated original-coordinate identities.
Now
\[
 F_{ij}=0,\qquad
 E_i=i\dot c_i\Phi+i c_i\dot\Phi,\qquad
 \sum_i[\mathcal A_i,E_i]
          =r^2\left(\sum_i c_i^2\right)\Omega.
\]
Thus the full Gauss residual, with precisely the convention in
\eqref{q444:full-gauss}, is
\[
 \boxed{\left(\frac{2r^2}{g^2}\sum_i c_i^2+r^2+B^2\right)\Omega.}
\]
Its coefficient is strictly positive and $\Omega\neq0$, proving the
obstruction even before the spatial gauge equations are used.

For completeness those equations force all $c_i$ to vanish as well:
here $J_i=F_{ij}=0$, so they require
$i(\ddot c_i\Phi+2\dot c_i\dot\Phi+c_i\ddot\Phi)=0$.
At a fixed time conjugate by the retained moving frame. The $(1,1)$
entries of $H$ and $[A,H]$ vanish, while that of $[A,[A,H]]$ is
$-2ax\neq0$. Hence $c_i=0$. The remaining Gauss residual is then
$(r^2+B^2)\Omega$, still nonzero.
\end{proof}

\begin{proposition}[Arbitrary temporal connection and exact Gauss inertia]
\label{q444:gauss-inertia}
For arbitrary homogeneous scalar and spatial-gauge positions and
velocities, define
\[
 J_{\mathrm{sc}}=
 [\Phi,\dot\Phi]+b\dot b^*-\dot b b^*,\qquad
 \mathscr I_{\Phi,b}Z=
 \bigl([\Phi,[\Phi,Z]]+bb^*Z+Zbb^*\bigr)_{\mathrm{tf}},
\]
\[
 \mathscr K_{\mathcal A}Z
       =-\sum_i[\mathcal A_i,[\mathcal A_i,Z]].
\]
Then Gauss' equation is exactly
\begin{equation}\label{q444:inertia-equation}
 \left(\mathscr I_{\Phi,b}+\frac2{g^2}\mathscr K_{\mathcal A}\right)
 \mathcal A_0
 =(J_{\mathrm{sc}})_{\mathrm{tf}}
       -\frac2{g^2}\sum_i[\mathcal A_i,\dot{\mathcal A}_i].
\end{equation}
The operator on the left is self-adjoint and positive semidefinite for
$\langle Z,W\rangle=-\operatorname{tr}(ZW)$ on $\mathfrak{su}(3)$.
Its kernel is exactly
\[
 \{Z:[Z,\Phi]=0,\ Zb=0,\ [Z,\mathcal A_i]=0\text{ for every }i\}.
\]
The right-hand side always annihilates that kernel. Thus Gauss' equation
has a pointwise solution for $\mathcal A_0$, unique modulo this kernel.
This pointwise statement asserts neither a smooth selection across
kernel-rank changes nor a solution of the remaining field equations.
\end{proposition}
\begin{proof}
Expansion gives
$J_0=J_{\mathrm{sc}}-\mathscr I_{\Phi,b}\mathcal A_0$
after trace-free projection, and
\[
 \sum_i[\mathcal A_i,E_i]
 =\sum_i[\mathcal A_i,\dot{\mathcal A}_i]
          +\mathscr K_{\mathcal A}\mathcal A_0.
\]
Substitution proves \eqref{q444:inertia-equation}. Cyclicity gives
\[
 \langle Z,\mathscr I_{\Phi,b}Z\rangle
       =\operatorname{tr}[Z,\Phi]^2+2|Zb|^2,\qquad
 \langle Z,\mathscr K_{\mathcal A}Z\rangle
       =\sum_i\|[\mathcal A_i,Z]\|_{\mathfrak{su}}^2,
\]
where $\|T\|_{\mathfrak{su}}^2=-\operatorname{tr}T^2$.
More explicitly, cyclicity gives the bilinear identities
\[
 \langle Z,\mathscr I_{\Phi,b}W\rangle
 =\operatorname{tr}([Z,\Phi][W,\Phi])
                  +2\operatorname{Re}((Zb)^*Wb),\qquad
 \langle Z,\mathscr K_{\mathcal A}W\rangle
 =\sum_i\langle[\mathcal A_i,Z],[\mathcal A_i,W]\rangle.
\]
These are symmetric in $Z,W$ and prove self-adjointness; the sums of
squares prove the stated kernel. If $Z$ is in the kernel, the matter moment
paired with $Z$ vanishes because $[Z,\Phi]=Zb=0$, and
$\operatorname{tr}(Z[\mathcal A_i,\dot{\mathcal A}_i])
=\operatorname{tr}([Z,\mathcal A_i]\dot{\mathcal A}_i)=0$.
Thus the right-hand side is orthogonal to the kernel. In a finite
dimensional real inner-product space a self-adjoint operator has image
equal to the orthogonal complement of its kernel: this follows by
diagonalizing its real symmetric matrix. It gives existence and the
claimed exact ambiguity.
\end{proof}

For the literal moving orbit, subtracting the original scalar equations
while retaining arbitrary $\mathcal A_0$ gives the exact additional
equations
\begin{align*}
 [\dot{\mathcal A}_0,\Phi]+2[\mathcal A_0,\dot\Phi]
 +[\mathcal A_0,[\mathcal A_0,\Phi]]
 -\sum_i[\mathcal A_i,[\mathcal A_i,\Phi]]&=0,\\
 \dot{\mathcal A}_0b+2\mathcal A_0\dot b+\mathcal A_0^2b
                   -\sum_i\mathcal A_i^2b&=0.
\end{align*}
Their norm consequences are
\[
 \sum_i\operatorname{tr}[\mathcal A_i,\Phi]^2
       =\operatorname{tr}X^2-\operatorname{tr}\dot\Phi^2,\qquad
 \sum_i|\mathcal A_i b|^2=|v|^2-|\dot b|^2.
\]
For example the trace pairing of the first equation uses
$\operatorname{tr}(\Phi[\dot{\mathcal A}_0,\Phi])=0$ and
$\operatorname{tr}(\Phi[\mathcal A_0,\dot\Phi])
=-\operatorname{tr}([\mathcal A_0,\Phi]\dot\Phi)$; expansion of $X^2$
then gives the first identity. The real pairing of the second equation
with $b$ gives the second by anti-Hermiticity. These right-hand sides
need not be zero; temporal gauge cannot be imposed while silently keeping
the same literal scalar functions.

Here is an exact counterexample to extending the temporal-gauge
stabilizer deduction to nonzero $\mathcal A_0$. For real constants
$\alpha,\beta_1,\beta_2,\beta_3$ satisfying
$\alpha^2-\sum_i\beta_i^2=1$, keep the same scalar orbit and set
\[
 \mathcal A_0=(\alpha-1)\Omega,\qquad
 \mathcal A_i=\beta_i\Omega.
\]
All curvatures vanish, while
$D_0\Phi=\alpha[\Omega,\Phi]$, $D_0b=\alpha\Omega b$,
$D_i\Phi=\beta_i[\Omega,\Phi]$, and $D_i b=\beta_i\Omega b$.
Their covariant wave operators therefore equal the original second
derivatives, proving both scalar equations exactly. Taking
$\alpha=5/3$, $\beta_1=4/3$, $\beta_2=\beta_3=0$ gives
$\mathcal A_1b\neq0$: the spatial field does not stabilize the scalars.
Nevertheless $J_0=\alpha J_{\mathrm{sc}}\neq0$, since $|\alpha|\ge1$,
and the zero electric field cannot satisfy Gauss' equation. These are
scalar-compatible flat connections, not full gauge solutions.

There is also no lift of the moving orbit with all spatial
$\mathcal A_i=0$, even when $\mathcal A_0$ is arbitrary.
Indeed $J_{\mathrm{sc}}=-\mathscr I_{\Phi,b}\Omega$, so
\eqref{q444:inertia-equation} gives
$\mathscr I_{\Phi,b}(\mathcal A_0+\Omega)=0$.
Its kernel description implies $D_0\Phi=D_0b=0$.
Equation~\eqref{q444:gauge-b} would then require
$(aI_3+\kappa\Phi)b=ab=0$, impossible because $a>0$ and $B>0$.
This argument, and Theorem~\ref{q444:temporal-no-lift}, do not exclude
the general case with both nonzero $\mathcal A_0$ and nonzero spatial
fields satisfying all coupled equations.

\subsection{Constructive homogeneous gauge solutions and their observations}

\begin{theorem}[A real radial family of full gauge solutions]
\label{q444:radial-gauge}
Let $D=\operatorname{diag}(-2,1,1)$ and retain
\[
 \Phi(t)=x(t)D,\qquad b(t)=\beta(t)e_1,\qquad
 x(t),\beta(t)\in\mathbb R,\qquad \mathcal A_\mu=0.
\]
This is a full homogeneous gauge--scalar solution exactly when
\begin{align*}
 \ddot x&=-a x-\kappa x^2+\frac{\kappa}{3}\beta^2,\\
 \ddot\beta&=-(a-2\kappa x)\beta,\qquad
 a=m^2+\Gamma(6x^2+2\beta^2).
\end{align*}
For every real initial position and velocity these equations have a
global solution. They give finite-energy homogeneous solutions on a
finite-volume flat $\mathbb T^3$, and local solutions on $\mathbb R^3$,
without asserting finite total energy on $\mathbb R^3$.
\end{theorem}
\begin{proof}
Here $[\Phi,\dot\Phi]=0$ and
$b\dot b^*-\dot b b^*=0$ because $x,\beta$ are real.
Thus every $J_\mu$ vanishes identically and so do all curvatures.
The gauge equations hold exactly, not after discarding Gauss' equation.
Directly
\[
 (\operatorname{adj}(xD)-\beta^2e_1e_1^*)_{\mathrm{tf}}
          =\left(-x^2+\frac{\beta^2}{3}\right)D,\qquad
 (aI_3+\kappa xD)\beta e_1=(a-2\kappa x)\beta e_1.
\]
The two scalar equations are consequently the displayed real system.
Its conserved density is
\[
 \mathcal E=3\dot x^2+\dot\beta^2+
 V(xD,\beta e_1),\qquad
 V=\frac{m^2}{2}s+\frac{\Gamma}{4}s^2
               -2\kappa x(\beta^2-x^2),\quad s=6x^2+2\beta^2.
\]
The stated conservation follows by differentiating and substituting the
two real equations; the $x$ kinetic coefficient is $3$, not $1/2$.
Moreover $|x|\le\sqrt{s/6}$ and $\beta^2+x^2\le s/2$, so
$|2x(\beta^2-x^2)|\le s^{3/2}/\sqrt6$ and
\[
 V\ge\frac{\Gamma}{4}s^2-\frac{|\kappa|}{\sqrt6}s^{3/2}.
\]
This is coercive and globally bounded below. Conservation bounds the
positions and velocities. The polynomial first-order vector field is
locally Lipschitz and bounded with bounded derivative on their compact
range; hence no finite-time escape is possible, proving global existence.
Multiplying the finite density by the finite volume of $\mathbb T^3$
gives finite total energy there.
\end{proof}

\begin{proposition}[A full-gauge vacuum and its marked Lorentz observation]
\label{q444:vacuum-observation}
At $\Gamma=m^2=1$, $\kappa=6$, let
\[
 x_0=\frac{3+\sqrt3}{12},\qquad
 \Phi=x_0\operatorname{diag}(-2,1,1),\qquad b=3x_0e_1.
\]
Together with $\mathcal A_\mu=0$ this is an equilibrium full gauge
solution at a global minimum of the scalar potential.
For $Q=\Theta(\Phi,b)$ its spectrum is $(4x_0,-2x_0,-2x_0)$.
Its spectral-floor matrix $H=Q+2x_0I_3$ has rank one and is positive.
The complex observation and hidden coordinate are
\[
 G=3x_0\begin{pmatrix}1&-1\\-1&2\end{pmatrix},
 \qquad \rho=3x_0,\qquad \det G=9x_0^2=|\rho|^2.
\]
More generally its $SU(3)$ orbit, parameterized by unit vectors
$e=ge_1$, has
\[
 \Phi=x_0(I_3-3ee^*),\quad b=3x_0e,\quad
 G=3x_0P(I_3-ee^*)P^*,\quad
 \det G=9x_0^2|u^*e|^2=|\rho|^2.
\]
For fixed $P$, this one gauge orbit contains both null and timelike
observations. The fixed-$P$ determinant is therefore not a
gauge-invariant mass observable.
\end{proposition}
\begin{proof}
For a nonzero real equilibrium with $\beta=3x>0$, the radial equations
reduce to
$24\Gamma x^2-2\kappa x+m^2=0$.
The displayed $x_0$ is the larger root at $(1,1,6)$.
Its quaternionic matrix is
\[
 Q=\begin{pmatrix}-2x_0&0&0\\0&x_0&3x_0\mathbf j\\
                       0&-3x_0\mathbf j&x_0\end{pmatrix}.
\]
The lower block has eigenvalues $x_0\pm3x_0$, proving the spectrum.
For an independent global-minimum check, every trace-free quaternionic
Hermitian matrix satisfies $54N(Q)^2\le s(Q)^3$: for eigenvalues summing
to zero, the exact identity is
$s^3-54N^2=2\prod_{i<j}(\lambda_i-\lambda_j)^2$.
Thus, with $r=\sqrt{s}\ge0$,
$V\ge r^2/2+r^4/4-2r^3/\sqrt6$, with equality at eigenvalues
$(2y,-y,-y)$, $y\ge0$, $r=\sqrt6y$.
On this equality family
$V=3y^2+9y^4-12y^3$ and
$dV/dy=6y(1-6y+6y^2)$.
The larger positive root is
$y=(3+\sqrt3)/6=2x_0$; the derivative is positive, negative,
then positive on the three successive intervals, and this minimum has
negative value, since at $y=2/3$ the potential is $-4/9$.
It is therefore the global minimum, attained by the displayed $Q$.

The complex part of $H$ is
$3x_0\operatorname{diag}(0,1,1)$, and compression by the retained $P$
gives the displayed $G$. The $SU(3)$ action on pairs is
$(\Phi,b)\mapsto(g\Phi g^*,gb)$; it is exactly quaternionic congruence
under $\Theta$, since $B(gb)=gB(b)g^{\mathsf T}$ for $g\in SU(3)$.
Consequently the spectrum and spectral-floor shift are unchanged.
Writing $v=Pe$ and using the determinant lemma in dimension two gives
\[
 \det(C-vv^*)=\det C\,(1-v^*C^{-1}v)
 =3\left(1-e^*\left(I_3-\frac{uu^*}{3}\right)e\right)
 =|u^*e|^2.
\]
This proves the general observation formula. A unit $e\in u^\perp$
gives determinant zero, while $e=e_1$ gives a positive determinant.
These vectors are related by $SU(3)$: extend either to a unitary basis
and adjust the phase of one of the remaining columns to make determinant
one. Thus the examples are genuinely on one gauge orbit.

The exact covariance map retains the observer as part of the data:
under $H'=gHg^*$, $P'=Pg^*$ one has
$P'H'P'^*=PHP^*$. Complex-component projection therefore has the same
value. Simultaneously $u'=gu$, $b'=gb$ gives $u'^*b'=u^*b$.
This is covariance of the marked observer, not invariance of a
fixed observer under an active transformation of the state.
\end{proof}

\begin{proposition}[Gauge averaging and distinct minimum orbits]
\label{q444:vacuum-average}
Let $d\mu$ be Haar measure on $SU(3)$ with its retained positive total
mass $V_\mu$, not set to one. For every complex Hermitian $A$,
\[
 \det(PAP^*)=u^*\operatorname{adj}(A)u,\qquad
 \int_{SU(3)}\det(PgAg^*P^*)\,d\mu(g)
                    =V_\mu\operatorname{tr}\operatorname{adj}(A).
\]
For a quaternionic positive rank-one matrix $H=\Theta(A,b)$,
$\operatorname{adj}(A)=bb^*$. Hence its averaged observation is
$V_\mu\|b\|^2$.

Put $\chi=2x_0>0$ for the preceding minimum and, for
$0\le\theta\le\pi/4$, define
\[
 v_\theta=(\cos\theta,\ \sin\theta\,\mathbf j,\ 0)^{\mathsf T},
 \qquad Q_\theta=\chi(3v_\theta v_\theta^*-I_3).
\]
Every $Q_\theta$ is a scalar global minimum, and its complex-plus-triplet
coordinate satisfies
\[
 b_\theta=(0,0,-3\chi\cos\theta\sin\theta)^{\mathsf T},\qquad
 \|b_\theta\|^2=\frac{9\chi^2}{4}\sin^2(2\theta).
\]
Consequently this minimum family contains pairwise $SU(3)$-inequivalent
states with invariant $\|b\|^2$ ranging from $0$ to $9\chi^2/4$.
The coefficient agrees with the supplied packet's invariant formula;
this calculation validates that coefficient rather than claiming it
as a newly discovered one.
\end{proposition}
\begin{proof}
The signed $2$-by-$2$ minors of $P$ form precisely $u=(1,1,1)^{\mathsf T}$.
The exterior-square form of the Cauchy--Binet expansion therefore gives
$\det(PAP^*)=u^*\operatorname{adj}(A)u$; equivalently this follows by
expanding the six entries of $PAP^*$ and collecting the cofactors.
The identity $\operatorname{adj}(gAg^*)=g\operatorname{adj}(A)g^*$
holds for invertible $A$ and then for every $A$ by polynomial continuity.
The Haar average of these conjugates commutes with every element of
$SU(3)$. Commutation with diagonal determinant-one phase matrices kills
its off-diagonal entries, and commutation with determinant-one cyclic
permutations makes its diagonal entries equal. Its trace fixes the
average at $V_\mu\operatorname{tr}(\operatorname{adj}A)I_3/3$.
Since $u^*u=3$, the stated integral follows with all factors retained.

Write $r=b^*b$. The skew convention gives directly
$B(b)\overline{B(b)}=bb^*-rI_3$.
The complex part of $H^2$ is therefore $A^2-bb^*+rI_3$,
and $e_2(H)=e_2(A)-r$.
For positive rank-one $H$, its spectral identity is
$H^2=(\operatorname{tr}H)H$ and $e_2(H)=0$.
Taking complex parts and using $\operatorname{tr}H=\operatorname{tr}A$
gives
$A^2-(\operatorname{tr}A)A+e_2(A)I_3=bb^*$.
The left side is $\operatorname{adj}(A)$, proving the assertion.

Finally $v_\theta^*v_\theta=1$, so $Q_\theta$ has fixed spectrum
$(2\chi,-\chi,-\chi)$ and hence the same minimum value.
Its $(1,2)$ quaternionic entry is
$-3\chi\cos\theta\sin\theta\,\mathbf j$, with the other quaternionic
off-diagonal coefficients zero. The retained definition of $B(b)$
therefore gives the displayed $b_\theta$ with its minus sign.
Taking its norm proves the formula. It is strictly increasing on the
given interval; since $b\mapsto gb$ preserves its norm, distinct
parameters there cannot lie on one $SU(3)$ orbit.
The positive spectral-floor matrix is $3\chi v_\theta v_\theta^*$,
so the averaging identity applies to the same triplet coordinate.
\end{proof}
